\begin{proof}[Proof of \cref{obj:lem:shifted-legendre-derivative}]
Fix \(j\ge1\). We first establish the derivative expansion, then bound it on the unit interval.

\begin{enumerate}
\item By the change of variables from \([0,1]\) to \([-1,1]\), the orthogonality in \cref{obj:lem:classical-legendre-facts} gives, for all nonnegative integers \(k,\ell\),
\[
\int_0^1 R_k(t)R_\ell(t)\,dt
=\frac12\int_{-1}^1
P_k^{\mathrm{Leg}}(t)P_\ell^{\mathrm{Leg}}(t)\,dt
=
\begin{cases}
\dfrac{1}{2k+1},&k=\ell,\\
0,&k\ne\ell.
\end{cases}
\]
Each polynomial \((-1)^kR_k\) has degree \(k\) and nonzero leading coefficient. Consequently,
\[
\operatorname{span}_{\mathbb R}\{(-1)^kR_k:0\le k<j\}
=
\{p\in\mathbb R[t]:\deg p<j\},
\]
where the set on the right includes the zero polynomial. Orthogonality and linearity therefore imply
\[
\int_0^1 p(t)R_j(t)\,dt=0
\qquad
\text{for every real polynomial \(p\) of degree below \(j\)}.
\]
Indeed, this integral vanishes on each spanning polynomial \((-1)^kR_k\), and hence on their linear combinations.
% lean: CausalSmith.Stat.RdTruesideNoiseFrontier.polynomial_shiftedLegendre_orthogonal

\item The same spanning argument gives the expansion
\[
p(t)=
\sum_{k=0}^{j-1}(2k+1)
\left(\int_0^1p(x)R_k(x)\,dx\right)R_k(t),
\qquad t\in\mathbb R,
\]
for every real polynomial \(p\) of degree below \(j\), including zero. To verify it on a spanning polynomial, take \(p=(-1)^\ell R_\ell\), where \(0\le\ell<j\). The right-hand side becomes
\[
\sum_{k=0}^{j-1}(2k+1)(-1)^\ell
\left(\int_0^1R_\ell(x)R_k(x)\,dx\right)R_k(t)
=(-1)^\ell R_\ell(t)=p(t).
\]
Both sides are linear in \(p\), so the identity extends to the entire span.
% lean: CausalSmith.Stat.RdTruesideNoiseFrontier.polynomial_shiftedLegendre_expansion

\item Polynomial differentiation gives
\[
\deg R_j'<j\quad(j\ge1),
\qquad R_0'=0.
\]
Applying the preceding expansion to \(R_j'\) yields
\[
R_j'(t)=
\sum_{k=0}^{j-1}(2k+1)
\left(\int_0^1R_j'(x)R_k(x)\,dx\right)R_k(t).
\]
For \(0\le k<j\), the polynomial \(R_k'\) has degree below \(j\), or is zero when \(k=0\). Thus the orthogonality established above supplies
\[
\int_0^1R_k'(x)R_j(x)\,dx=0.
\]
Integration by parts now gives
\[
\int_0^1R_j'(x)R_k(x)\,dx
=
R_k(1)R_j(1)-R_k(0)R_j(0)
-\int_0^1R_k'(x)R_j(x)\,dx.
\]
The endpoint identity and reflection symmetry in
\cref{obj:lem:classical-legendre-facts} imply, for every \(k\ge0\),
\[
R_k(0)=P_k^{\mathrm{Leg}}(-1)=(-1)^k,
\qquad
R_k(1)=P_k^{\mathrm{Leg}}(1)
=(-1)^kP_k^{\mathrm{Leg}}(-1)=1.
\]
Since \(j=(j-k)+k\), the endpoint product satisfies
\[
(-1)^k(-1)^j=(-1)^{j-k}(-1)^{2k}=(-1)^{j-k}.
\]
Consequently,
\[
\int_0^1R_j'(x)R_k(x)\,dx
=1-(-1)^{j-k}
=
\begin{cases}
2,&j-k\text{ odd},\\
0,&j-k\text{ even}.
\end{cases}
\]
Substituting these coefficients into the expansion proves, for every real \(t\),
\[
R_j'(t)=
2\sum_{\substack{0\le k<j\\j-k\ \mathrm{odd}}}(2k+1)R_k(t).
\]
% lean: CausalSmith.Stat.RdTruesideNoiseFrontier.shiftedLegendre_derivative_expansion

\item Fix \(t\in[0,1]\). Then \(2t-1\in[-1,1]\), so the unit bound in \cref{obj:lem:classical-legendre-facts} gives
\[
|R_k(t)|=|P_k^{\mathrm{Leg}}(2t-1)|\le1
\qquad(k\ge0).
\]
Using the expansion, the triangle inequality, and \(2k+1>0\), we obtain
\[
\begin{aligned}
|R_j'(t)|
&\le
2\sum_{\substack{0\le k<j\\j-k\ \mathrm{odd}}}
|(2k+1)R_k(t)|\\
&=
2\sum_{\substack{0\le k<j\\j-k\ \mathrm{odd}}}
(2k+1)|R_k(t)|\\
&\le
2\sum_{\substack{0\le k<j\\j-k\ \mathrm{odd}}}(2k+1).
\end{aligned}
\]
% lean: CausalSmith.Stat.RdTruesideNoiseFrontier.shiftedLegendre_deriv_bound_of_expansion

\item To evaluate the last sum, define, for every integer \(\ell\ge0\),
\[
\omega_\ell:=
\sum_{\substack{0\le k<\ell\\\ell-k\ \mathrm{odd}}}(2k+1).
\]
First,
\[
\sum_{k=0}^{\ell-1}(2k+1)=\ell^2
\qquad(\ell\ge0).
\]
This follows by induction: the empty sum at \(\ell=0\) is zero, and the induction step is
\[
\sum_{k=0}^{\ell}(2k+1)
=\ell^2+(2\ell+1)=(\ell+1)^2.
\]
For each \(k<\ell\), the integers \(\ell-k\) and \(\ell+1-k\) have opposite parity. The additional term \(k=\ell\) belongs to \(\omega_{\ell+1}\). Hence
\[
\omega_{\ell+1}+\omega_\ell
=
\sum_{k=0}^{\ell-1}(2k+1)+(2\ell+1)
=(\ell+1)^2.
\]
We now prove \(2\omega_\ell=\ell(\ell+1)\) by induction. The base case follows from \(\omega_0=0\); if the identity holds at \(\ell\), then
\[
2\omega_{\ell+1}
=2(\ell+1)^2-2\omega_\ell
=2(\ell+1)^2-\ell(\ell+1)
=(\ell+1)(\ell+2).
\]
In particular, \(2\omega_j=j(j+1)\). Substitution into the preceding bound proves
\[
|R_j'(t)|\le j(j+1)
\qquad(t\in[0,1]).
\]
% lean: CausalSmith.Stat.RdTruesideNoiseFrontier.shiftedLegendre_parity_weight_sum
\end{enumerate}
\end{proof}
